12 \hfill 1. A first look at Banach and Hilbert spaces

Now choosing $h=0$ (and using $s(f,0) = 0$) shows $s(f,g) = 2s(f,\frac{g}{2})$ and
thus $s(f,g) + s(f,h) = s(f,g+h)$. Furthermore, by induction we infer
$\lambda s(f,g) = s(f, \lambda g)$ for every positive rational $\lambda$.
By continuity (check this!) this holds for all $\lambda > 0$ and $s(f,-g) = -s(f,g)$
respectively $s(f, ig) = is(f,g)$ finishes the proof. $\square$

Note that the parallelogram law and the polarization identity even hold
for skew linear forms.

But how do we define a scalar product on $\mathcal{C}(I)$? One possibility is
\[
\langle f,g \rangle = \int_a^b f^*(x)g(x)dx.
\tag{1.32}
\]
The corresponding inner product space is denoted by $L^2(I)$. Note that we
have
\[
\|f\|_\infty \leq \sqrt{b-a} \|f\|_2
\tag{1.33}
\]
and hence the maximum norm is stronger than the $L^2$ norm.

Suppose we have two norms $\|\cdot\|_1$ and $\|\cdot\|_2$ on a space $\mathcal{X}$. Then $\|\cdot\|_2$ is
said to be \textbf{stronger} than $\|\cdot\|_1$ if there is a constant $m > 0$ such that
\[
\|f\|_1 \leq m \|f\|_2.
\tag{1.34}
\]
It is straightforward to check that

\noindent\textbf{Lemma 1.7.} If $\|\cdot\|_2$ is stronger than $\|\cdot\|_1$, then any $\|\cdot\|_2$ Cauchy sequence
is also a $\|\cdot\|_1$ Cauchy sequence.

Hence if a function $F: \mathcal{X} \to \mathcal{Y}$ is continuous in $(\mathcal{X}, \|\cdot\|_1)$ it is also
continuos in $(\mathcal{X}, \|\cdot\|_2)$ and if a set is dense in $(\mathcal{X}, \|\cdot\|_2)$ it is also dense in
$(\mathcal{X}, \|\cdot\|_1)$.

In particular, $L^2$ is separable. But is it also complete? Unfortunately
the answer is no:

\noindent\textbf{Example.} Take $I = [0, 2]$ and define
\[
f_n(x) = \begin{cases}
0, & 0<x<1-\frac{1}{n} \\
1 + n(x-1), & 1-\frac{1}{n}\leq x \leq 1 \\
1, & 1 < x \leq 2
\end{cases}
\tag{1.35}
\]
then $f_n(x)$ is a Cauchy sequence in $L^2$, but there is no limit in $L^2$! Clearly
the limit should be the step function which is $0$ for $0 < x < 1$ and $1$ for
$1 \leq x \leq 2$, but this step function is discontinuous (Problem 1.8)! $\diamondsuit$

This shows that in infinite dimensional spaces different norms will give
raise to different convergent sequences! In fact, the key to solving prob-
lems in infinite dimensional spaces is often finding the right norm! This is
something which cannot happen in the finite dimensional case.